\section{What the proof makes reusable}
\label{p1:discussion}

The exponent \(4/3\) results from the relation between two scales:
the leading pair mass is \(N^{-2/3+o(1)}\), and the quadratic
reference omits a contribution formed from two such interactions.
The arithmetic work shows that all other terms relevant to the
reconstruction can be controlled at this precision. The proof also
separates several choices that can be made in other rare-event
problems.

\subsection{Choose a reference before estimating the error}

The independent Bernoulli reference retains every exact marginal.
Its centered expansion has no one-site term, and repeated indices
are already accounted for. The retained pair correction then
removes the complete second-order contribution. This makes the
quantity to be bounded a residual interaction, with a specified
sign and normalization.

A homogeneous envelope provides a second useful reduction. The
joint probability of the homogeneous constraints bounds the actual
joint probability, the exact marginal rank weight and the absolute
centered activity. Geometric counting can therefore be performed
once and used for each of these weights. At a fixed arity \(r\),
the comparison with the nominal weight costs at most \(2^r\).
That fixed factor can be absorbed into the constant of a fixed-order
estimate. The estimates that sum over unbounded arities keep their
own weights and quantifiers.

Resultants and degree-uniform curve bounds also appear in Shparlinski
and Sleiman's count of consecutive multiplicatively dependent triples
\cite[Theorem~1.1 and Section~3]{shparlinskisleiman2026}.
The additional feature used here is the simultaneous placement of a
large space of square-class relations on one common curve, together
with the integer fibres that those relations force in the offset grid.
This converts geometric information about logarithmic windows into a
bound for their shared rank.

\subsection{Stop where the remaining dependence is still positive}

The specialization of the abstract lace expansion
\cite{zeilberger1997} in Section~\ref{p1:stop:section} retains a
hereditary class of subsets with computable joint probabilities.
At the first excluded prefix, all later variables remain in a
positive conditional measure. Its total mass is the joint
probability of the prefix. The boundary therefore has a direct
probabilistic weight even though the retained expansion is signed.
Incompatible pair and triple terms must still be combined before
absolute values are taken.

To use this identity in another model, one needs a retained class
that matches the chosen reference and quantitative control of the
first excluded configurations. Here compression makes a dependent
pair one affine block. Conditional arithmetic counts control a
fixed number of blocks, and congruence certificates control the
large-cardinality part. The reconstruction estimate records exactly
what these inputs accomplish:
\[
 \|\Law(Z_N)-Q_{\mu_N,\alpha_N}\|_1
       \ll N^\varepsilon A_3+N^{-4/3+\varepsilon}.
\]
The finite identity itself does not give the arithmetic bounds
needed on its right-hand side.

\subsection{Separate fixed orders from a uniform tail}

A fixed-order theorem allows its constant and scale threshold to
depend on the order. It cannot be summed up to an order growing
with \(N\) without further information. The proof instead chooses
a finite cutoff after the desired accuracy, treats the finitely
many lower orders by their individual estimates, and uses one
uniform weighted tail above that cutoff.

Proposition~\ref{p1:uniform:full-support} illustrates this
arrangement. Finite layers give a half power for each effective
block, while the independent tail gives every power strictly below
\(2\). Their combination yields the exponent
\(\min(s/2,2)\), with an arbitrary fixed margin and a common
constant for all \(s\ge2\). This is a bound on a deterministic
sum of independent counting weights. It is not an exponential-tail
statement for the law of \(Z_N\).

The independent fixed-moment proof has a different use: it turns a
polynomial test into a norm lower bound and gives fixed polynomially
weighted consequences of the main theorem. The finite-difference
example in Section~\ref{p1:moment-obstruction} shows why even very strong information at
each fixed moment order does not by itself provide a polynomial
rate for the complete law.

\subsection{Test the approximation that was actually retained}

The fourth Charlier polynomial detects
\(\kappa_{4,N}+3\alpha_N^2\) after subtraction of the quadratic
reference. The nonnegative compatible part of the quartic identity
prevents cancellation at the scale of \(\alpha_N^2\).
For the exact Hermite reference the same test instead detects
\(\kappa_{4,N}\), because its exponential already contains the
two-pair term. A lower bound is therefore tied to a specified
reference, not just to the underlying dependent count.

The full equivalent for the exact quadratic remainder, the full
constant for the third factorial cumulant, and the sharp error of
the exact Hermite reference remain open in this work. So do the
errors after cubic or quartic correction and the optimization of
total variation over all Poisson means. The explicit channel
coefficients determine reference laws, but the constants in the
asymptotic remainders and the thresholds in \(N\) are not supplied
numerically. Applying the method elsewhere requires counterparts
of the conditional counts, compatibility cancellations and uniform
boundary estimate; the stopping identity alone does not impose a
universal approximation exponent.
